% META: {"id": "TSPE-GEOESPACE-EX-040", "chapitre": "TSPE-GEOMETRIE-ESPACE", "type_objet": "exercice", "capacites": ["TSPE-GEOMETRIE-ESPACE-C13"], "capacites_codes": ["C13"], "methodes": ["M1"], "parcours": 2, "competences": ["raisonner", "communiquer"], "duree_min": 8, "mode_creation": "ex_nihilo", "sources_inspiration": [], "fichier_tex": "chapitres/TSPE-GEOMETRIE-ESPACE/exercices/TSPE-GEOESPACE-EX-040.tex", "corrige_tex": "chapitres/TSPE-GEOMETRIE-ESPACE/corriges/TSPE-GEOESPACE-CO-040.tex", "status": "approved"}
% BEGIN-VERIFY
% assert True
% END-VERIFY
\begin{exercice}{TSPE-GEOESPACE-EX-005}{3}{16}
Soit $\mathcal{P}$ un plan, $M$ un point n'appartenant pas a $\mathcal{P}$, et $H$ le projete orthogonal de $M$ sur $\mathcal{P}$. Soit $N$ un point de $\mathcal{P}$ différent de $H$.
\begin{enumerate}
  \item Justifier que $\vec{HN}$ est un vecteur du plan $\mathcal{P}$.
  \item Justifier que $\vec{MH}\cdot\vec{HN}=0$.
  \item En appliquant le théorème de Pythagore dans le triangle $MHN$, exprimer $MN^2$ en fonction de $MH^2$ et $HN^2$.
  \item En déduire que $MN>MH$.
  \item Conclure : que représente ce résultat pour le point $H$ ?
\end{enumerate}
\end{exercice}
